\chapter{Price Theory at Work}\label{ch:b1:price-theory-at-work}

In 2020 a flat in France cost about forty years of its rent; in 1950 it had
cost under fifteen. In the United States the number was about twenty-four
at both dates. % figdata: rent_points
Why does a dwelling cost a given number of years of rent, and why did that
number move so much in some countries and not at all in others? The answer
is an \omterm{def:b1:price-theory-at-work:arbitrage}{arbitrage}: buying and renting are two ways of getting the same thing,
and their prices cannot drift apart for long. This chapter puts the price
theory of the last four chapters to work on such questions: what a price
does when the law stops it moving, how a quota differs from a tariff, why
a stored crop gets dearer month by month, and what fixes the price of a
house.

\begin{openingchart}{Housing rental yield and the real long-term interest
rate, France, United Kingdom and United States, 1950--2020; the real rate is
the long-term government bond rate minus consumer-price inflation averaged
over three years.}{Jordà--Schularick--Taylor Macrohistory Database, release 6}
\begin{tikzpicture}
\begin{groupplot}[group style={group size=2 by 1, horizontal sep=16mm}, omtime, xmin=1950, xmax=2020,
  width=0.46\linewidth, height=0.4\linewidth, xtick={1950,1970,1990,2010}]
\nextgroupplot[ylabel={Rental yield (\%)}, ymin=0, ymax=8, title={\small Rent as a share of price}, xlabel={Year},
  width=0.46\linewidth, legend style={font=\footnotesize, at={(1.1,-0.28)}, anchor=north, legend columns=3}]
  \addplot+[no marks, thick] table[x=year, y=yield_fra, col sep=comma] {figdata/bachelor-1/14-price-theory-at-work/out/rent_yield.csv};
  \addplot+[no marks, thick] table[x=year, y=yield_gbr, col sep=comma] {figdata/bachelor-1/14-price-theory-at-work/out/rent_yield.csv};
  \addplot+[no marks, thick] table[x=year, y=yield_usa, col sep=comma] {figdata/bachelor-1/14-price-theory-at-work/out/rent_yield.csv};
  \legend{France, United Kingdom, United States}
\nextgroupplot[ylabel={Real long rate (\%)}, ymin=-8, ymax=8, title={\small Real interest rate}, xlabel={Year},
  width=0.46\linewidth, restrict y to domain=-8:8, unbounded coords=jump]
  \draw[gray, dotted] (axis cs:1950,0) -- (axis cs:2020,0);
  \addplot+[no marks, thick] table[x=year, y=real_fra, col sep=comma] {figdata/bachelor-1/14-price-theory-at-work/out/rent_yield.csv};
  \addplot+[no marks, thick] table[x=year, y=real_gbr, col sep=comma] {figdata/bachelor-1/14-price-theory-at-work/out/rent_yield.csv};
  \addplot+[no marks, thick] table[x=year, y=real_usa, col sep=comma] {figdata/bachelor-1/14-price-theory-at-work/out/rent_yield.csv};
\end{groupplot}
\end{tikzpicture}
\end{openingchart}

The French yield fell from $6.79\%$ in 1950 to $2.49\%$ in 2020, the British
from $3.68\%$ to $2.22\%$, while the American one ended where it began, at
$4.23\%$. In 2000 the three real rates were between $3.69\%$ and $4.44\%$; in 2020
all three were below zero. % figdata: rent_points
France's real rate was far below the chart in 1950, at $-19.26\%$, after the
inflation of the war years. % figdata: rent_points
The chart is a test of the \omterm{def:b1:price-theory-at-work:arbitrage}{arbitrage} of \cref{sec:b1:price-theory-at-work:buy-rent}:
yields should follow the real rate, and they did so only in part.

\begin{recall}
Price ceilings and floors, and \omterm{def:b1:price-theory-at-work:arbitrage}{arbitrage} and storage as what prices do,
were met at school; present value at school too. The \omterm{def:b1:competitive-industry:capitalisation}{capitalisation} of a
rent into an asset price is \cref{def:b1:competitive-industry:capitalisation};
surplus and the \omterm{def:b1:partial-equilibrium-welfare:dwl}{deadweight loss} are in \cref{ch:b1:partial-equilibrium-welfare}.
\end{recall}

\section{The full price}\label{sec:b1:price-theory-at-work:full-price}

What a buyer gives up for a good is rarely only money.

\begin{definition}[Full price]\label{def:b1:price-theory-at-work:full-price}
\index{full price}
The \emph{full price} of a good is its money price plus the money value of
everything else the buyer must give up to get it: time spent queuing or
travelling, at the buyer's value of time, and any loss of quality.
\end{definition}

Demand responds to the \omterm{def:b1:price-theory-at-work:full-price}{full price}. A museum with free admission and a queue of
two hours is not free to someone whose time is worth 15 dalens an hour: it
costs 30 dalens. A free road that is congested costs its users time, and a
toll road beside it fills until the toll plus its shorter journey costs
about as much as the free road's longer one. When the law holds the money
price down, something else --- the queue, the quality, a side payment ---
rises until the \omterm{def:b1:price-theory-at-work:full-price}{full price} clears the market.

\section{Around a controlled price}\label{sec:b1:price-theory-at-work:controls}

A ceiling below the \omterm{def:b1:models-and-data:equilibrium}{equilibrium} price makes the quantity demanded exceed the
quantity supplied. Sellers supply only $S(\bar p)$, and the units are rationed
by something other than the money price.

\begin{omfigure}
\begin{tikzpicture}
\begin{axis}[omrecord, grid=none, xlabel={Quantity}, ylabel={Price}, xmin=0, xmax=80, ymin=0, ymax=80,
  width=0.52\linewidth, height=0.42\linewidth]
  \fill[omProp!35] (axis cs:20,60) -- (axis cs:40,40) -- (axis cs:20,20) -- cycle;
  \fill[omDef!18] (axis cs:0,20) rectangle (axis cs:20,60);
  \addplot[demand, domain=0:80] {80 - x};
  \addplot[supply, domain=0:80] {x};
  \addplot[guide] coordinates {(0,20) (60,20)};
  \addplot[guide] coordinates {(20,0) (20,60) (0,60)};
  \node[font=\footnotesize, anchor=west] at (axis cs:61,20) {ceiling};
  \node[font=\footnotesize] at (axis cs:10,40) {queue};
  \node[font=\footnotesize, anchor=west] at (axis cs:33,52) {DWL};
  \node[font=\footnotesize, anchor=south] at (axis cs:40,20) {shortage};
\end{axis}
\end{tikzpicture}

{\small A ceiling of 20 with demand $Q = 80 - p$ (blue) and supply $Q = p$ (red).
Sellers offer 20 units and buyers want 60. The last buyer served values a
unit at 60, so queues grow until a unit costs 60 in money and time: the
blue rectangle, worth 800, is time wasted in the queue, on top of the
triangle of lost trades, 400.}
\end{omfigure}

Who gains? Buyers who get the good at the ceiling without queuing; who
they are depends on the rationing: first come, by lottery, by connections.
Sellers lose. If rationing is by queuing the rectangle of
\cref{exo:b1:price-theory-at-work:ceiling} is lost too, and if units go to
buyers at random rather than to those who value them most, the loss is
larger still (\cref{exo:b1:price-theory-at-work:random}).

\begin{definition}[Parallel market]\label{def:b1:price-theory-at-work:parallel-market}
\index{parallel market}\index{black market|see{parallel market}}
A \emph{parallel market} is an unrecorded market in which a good under a
price ceiling is sold above the ceiling.
\end{definition}

Over time, sellers respond to a ceiling on more margins than quantity. A
landlord whose rent is capped spends less on upkeep, so the quality of the
dwelling falls towards what the rent pays for. In the long run, when supply
is more elastic, the shortage grows: owners can sell flats to owner
occupiers, convert them, or build none.

\begin{record}[Rent control in San Francisco]
In 1994 San Francisco extended rent control to small multi-family buildings
built before 1980, giving economists a comparison of similar buildings on
either side of the line. Rent control cut tenants' moves by about $20\%$ ---
the protection it was meant to give. But landlords whose buildings came
under control reduced the supply of rental housing on them by $15\%$, by
selling to owner occupiers and by redeveloping (Diamond, McQuade and Qian
2019). % ledger: b1:diamond2019
The model's long-run prediction, a fall in supply, was found; the margin by
which it happened was conversion rather than neglect.
\end{record}

A price floor --- a minimum price for a farm product, say --- has the
mirror effects: a surplus that someone must buy or destroy, and effort
spent by sellers competing for the buyers there are.

\begin{model}{A price ceiling}
\assumes A competitive market and a binding ceiling enforced on the money
price.
\predicts A shortage; rationing by queues, connections or side payments;
lower quality; a \omterm{def:b1:price-theory-at-work:parallel-market}{parallel market}; and a larger fall in supply in the long
run than in the short.
\tested Extending rent control in San Francisco cut the supply of rental
housing in the buildings concerned by $15\%$ (Diamond, McQuade and Qian
2019). % ledger: b1:diamond2019
\breaks Where owners cannot exit, convert or let quality slip, the
short-run supply response is small; where the ceiling is above the market
price, nothing happens at all.
\end{model}

\section{Quota against tariff}\label{sec:b1:price-theory-at-work:quota}

A country can limit imports of a good with a tariff, a tax on each unit
imported, or with a quota, a cap on the number. For any tariff there is a
quota that gives the same imports, the same domestic price and the same
quantities: the one equal to the imports under the tariff. The two differ
in who gets the gap between the world price and the domestic price.

\begin{definition}[Quota rent]\label{def:b1:price-theory-at-work:quota-rent}
\index{quota rent}
The \emph{quota rent} is the difference between the domestic and the
world price, times the quantity imported under a quota: the income of
whoever holds the right to import.
\end{definition}

Under a tariff the same rectangle is revenue to the state. Under a quota
it goes to the licence holders: domestic importers if the licences are
given to them, the state if it auctions them, or foreign exporters if the
exporting country runs the quota.

\begin{omfigure}
\begin{tikzpicture}
\begin{axis}[omrecord, grid=none, xlabel={Quantity}, ylabel={Price}, xmin=0, xmax=80, ymin=0, ymax=24,
  width=0.52\linewidth, height=0.42\linewidth]
  \fill[omDef!22] (axis cs:24,10) rectangle (axis cs:60,12);
  \addplot[demand, domain=0:80] {(120 - x)/5};
  \addplot[supply, domain=0:48] {x/2};
  \addplot[guide] coordinates {(0,10) (70,10)};
  \addplot[guide] coordinates {(0,12) (60,12)};
  \addplot[guide] coordinates {(24,0) (24,12)};
  \addplot[guide] coordinates {(60,0) (60,12)};
  \node[font=\footnotesize, anchor=north west] at (axis cs:1,9.8) {world price};
  \node[font=\footnotesize, anchor=north] at (axis cs:42,9.6) {revenue or quota rent};
\end{axis}
\end{tikzpicture}

{\small Demand $Q = 120 - 5p$ (blue), domestic supply $Q = 2p$ (red), world price 10.
A tariff of 2 or a quota of 36 units gives the same price, 12, and the same
imports, 36; the rectangle, 72, is tariff revenue in one case and \omterm{def:b1:price-theory-at-work:quota-rent}{quota
rent} in the other.}
\end{omfigure}

\begin{record}[A voluntary export restraint]
From 1981 Japanese carmakers limited their exports to the United States
under an agreement with the US government, a quota run by the exporting
country. A study of its effects from 1981 to 1990 measured the cost to
American buyers, the profits of the firms, and the tariff revenue the
United States gave up by choosing a quota over a tariff (Berry,
Levinsohn and Pakes 1999). % ledger: b1:berry1999
Under such a restraint the \omterm{def:b1:price-theory-at-work:quota-rent}{quota rent} goes to the exporting firms: they
can sell the limited number of cars at the high domestic price.
\end{record}

\section{Arbitrage over space and over time}\label{sec:b1:price-theory-at-work:arbitrage}

\begin{definition}[Arbitrage condition]\label{def:b1:price-theory-at-work:arbitrage}
\index{arbitrage condition}\index{arbitrage}
A trader who can buy a good in market $A$ and sell it in market $B$ at a
cost $c$ a unit earns a profit whenever $p_B > p_A + c$. The \emph{arbitrage
condition} is $p_B \le p_A + c$, with equality whenever the good actually flows
from $A$ to $B$. Over time, with $B$ a later date, $c$ is the cost of storing
the good and the interest on its value.
\end{definition}

Two regions linked by trade thus have prices that differ by the transport
cost, or less if no good flows. If Brenland's demand for imports is
$2(220 - p_B)$ and Kestria's export supply is $2(p_K - 160)$, a transport cost
of 10 dalens gives $p_K = 185$, $p_B = 195$ and a trade of 50 units; at a cost of
60 or more, nothing moves (\cref{pb:b1:price-theory-at-work:store}).

\begin{definition}[Cost of carry and convenience yield]\label{def:b1:price-theory-at-work:carry}
\index{cost of carry}\index{convenience yield}
The \emph{cost of carry} of a storable good from one date to the next is the
interest on its value plus the cost of storing it. The \emph{convenience
yield} is the benefit of holding a stock on hand --- not running out, being
able to sell at once --- measured as the amount by which the price rises
less than the cost of carry while stocks are held.
\end{definition}

With storage, a crop harvested once a year becomes dearer month by month:
if stocks are carried from September to July, the price must rise each
month by the \omterm{def:b1:price-theory-at-work:carry}{cost of carry}, net of the \omterm{def:b1:price-theory-at-work:carry}{convenience yield}, or no one would
hold them. At the next harvest the price falls back. The price of a futures
contract for delivery at a later date obeys the same condition; the
quantitative finance series studies it.

\begin{omfigure}
\begin{tikzpicture}
\begin{axis}[omrecord, grid=none, xlabel={Months from the first harvest}, ylabel={Dalens a tonne}, xmin=0, xmax=36,
  ymin=190, ymax=240, width=0.62\linewidth, height=0.36\linewidth, xtick={0,12,24,36}]
  \addplot[demand, samples=12, domain=0:11] {200*1.005^x + 2*(1.005^x - 1)/0.005};
  \addplot[demand, samples=12, domain=12:23] {200*1.005^(x-12) + 2*(1.005^(x-12) - 1)/0.005};
  \addplot[demand, samples=12, domain=24:35] {200*1.005^(x-24) + 2*(1.005^(x-24) - 1)/0.005};
  \addplot[guide] coordinates {(11,233.84) (12,200)};
  \addplot[guide] coordinates {(23,233.84) (24,200)};
\end{axis}
\end{tikzpicture}

{\small The price of a stored grain under the \omterm{def:b1:price-theory-at-work:arbitrage}{arbitrage condition} (the grain
store of \cref{pb:b1:price-theory-at-work:store}): 200 dalens a tonne at
harvest, then a rise each month by interest at $0.5\%$ and storage at 2 dalens,
to 233.84 eleven months later; the harvest brings it back to 200.}
\end{omfigure}

Retail prices show the pattern, and its absence where storage is
impossible.

\begin{omchart}{Seasonal pattern of three retail food prices, United States,
2000--2019: average deviation of each month's price from its year's mean
(log points).}{US Bureau of Labor Statistics, Average Price Data (series
AP), US city average}
\begin{tikzpicture}
\begin{axis}[omtime, xlabel={Month}, ylabel={Deviation from the year's mean (\%)}, xmin=1, xmax=12,
  ymin=-25, ymax=40, width=0.7\linewidth, height=0.4\linewidth, xtick={1,...,12}, legend pos=north west,
  legend style={font=\footnotesize}]
  \draw[gray, dotted] (axis cs:1,0) -- (axis cs:12,0);
  \addplot+[thick] table[x=month, y=strawberries, col sep=comma] {figdata/bachelor-1/14-price-theory-at-work/out/seasonal.csv};
  \addplot+[thick] table[x=month, y=potatoes, col sep=comma] {figdata/bachelor-1/14-price-theory-at-work/out/seasonal.csv};
  \addplot+[thick] table[x=month, y=rice, col sep=comma] {figdata/bachelor-1/14-price-theory-at-work/out/seasonal.csv};
  \legend{Strawberries, White potatoes, Long-grain rice}
\end{axis}
\end{tikzpicture}
\end{omchart}

Strawberries cannot be stored: their price follows the season of supply,
lowest in June and highest in December, a range of 53 points. Potatoes are
stored through the year: their price rises from a low in February to a
high in August, just before the main harvest, a range of 8.4 points.
Rice is stored for years and its price barely has a season: 3.6 points.
% figdata: seasonal_range
The potato path is not the saw-tooth of the figure above: the price keeps
falling for some months after the autumn harvest, as the year's crop
arrives in shops. But from February to August it rises by about $1.4$ points
a month, $8.4$ points over six months, the order of a \omterm{def:b1:price-theory-at-work:carry}{cost of carry}.
% figdata: seasonal_range

\section{Buying against renting}\label{sec:b1:price-theory-at-work:buy-rent}

A dwelling can be rented for $R$ a year or bought for $P$. Owning for a year
costs the interest given up on $P$, at the real rate $r$, plus depreciation
and upkeep $\delta P$, minus the expected real rise in the price, $gP$.

\begin{definition}[Rental yield and the rent--price arbitrage condition]\label{def:b1:price-theory-at-work:rental-yield}
\index{rental yield}\index{rent--price arbitrage condition}
The \emph{rental yield} of a dwelling is its yearly rent divided by its
price, $R/P$. The \emph{rent--price \omterm{def:b1:price-theory-at-work:arbitrage}{arbitrage condition}} says that renting and
owning cost the same:
\[
R = P\,(r + \delta - g), \qquad \frac RP = r + \delta - g.
\]
\end{definition}

If owning were cheaper, renters would buy and the price would rise; if
renting were cheaper, owners would sell and rent. The condition is the
\omterm{def:b1:competitive-industry:capitalisation}{capitalisation} of \cref{def:b1:competitive-industry:capitalisation} with
depreciation and expected growth added. It predicts that the price is
$R/(r + \delta - g)$ years of rent, and so that a fall in the real rate raises
the price of a given rent: with $\delta = 2\%$ and $g = 0$, a fall in $r$ from $4\%$ to
$2\%$ takes the price from 16.67 to 25 years of rent, a rise of $50\%$.
% computed: code/bachelor-1/14-price-theory-at-work (exercise rate-fall)

\begin{measure}[A rental yield]
The yields of the opening chart come from a dataset of total returns on
housing, equities, bonds and bills in 16 advanced economies since 1870
(Jordà et al. 2019). % ledger: b1:jorda2019
A yield there is a rent net of costs divided by the price of the
dwelling; long series are built by linking the ratio in a year when it is
well measured to indices of rents and of house prices. Errors in either
index accumulate along the series, so the level of a yield decades away
from its benchmark is less certain than its movements.
\end{measure}

The opening chart tests the condition. Over 1950--2020 the correlation of
the yield with the real rate is $0.63$ in the United States, $0.12$ in the
United Kingdom and $-0.15$ in France. % figdata: rent_corr
In the United States yields rose and fell with the real rate; in France and
Britain the yield fell through decades when the real rate did not, perhaps
because expected growth $g$ rose with years of rising prices, a term the
chart cannot see, or because the long yield series drift (see the measure
box).

\begin{model}{Arbitrage}
\assumes Traders who can buy, store, ship or rent at known costs, and who
act on any price gap larger than those costs.
\predicts Price gaps between places equal to transport costs where goods
flow; prices of stored goods rising through the crop year by the \omterm{def:b1:price-theory-at-work:carry}{cost of
carry}; \omterm{def:b1:price-theory-at-work:rental-yield}{rental yields} equal to the real rate plus depreciation, less
expected growth.
\tested Stored foods have small seasonal price swings and perishable ones
large swings (the chart above); \omterm{def:b1:price-theory-at-work:rental-yield}{rental yields} followed the real rate in
the United States from 1950 to 2020 but not in France or Britain (opening
chart). Price gaps across countries and distance are studied with the
gravity model of Year~3.
\breaks When shorting or storage is impossible --- a perishable crop, a
house that cannot be sold short --- and when expectations of price growth
swing, as in housing in the 1980s and 2000s.
\end{model}

\section*{Exercises}\label{sec:b1:price-theory-at-work:exercises}

\begin{exercise}[$\star$]\label{exo:b1:price-theory-at-work:full}
A concert ticket costs 4 dalens but requires two hours in a queue. What is
its \omterm{def:b1:price-theory-at-work:full-price}{full price} to someone whose time is worth 15 dalens an hour?
\end{exercise}

\begin{exercise}[$\star$]\label{exo:b1:price-theory-at-work:read}
From the opening chart's table: France's yield was $3.93\%$ in 1980 and
$2.49\%$ in 2020; its real rate was $1.97\%$ and $-1.29\%$. % figdata: rent_points
Compute the price in years of rent at both dates and the change in the
real rate.
\end{exercise}

\begin{exercise}[$\star$]\label{exo:b1:price-theory-at-work:quota-rent}
A quota allows 40 units at a world price of 10; the domestic price is 14.
Compute the \omterm{def:b1:price-theory-at-work:quota-rent}{quota rent}.
\end{exercise}

\begin{exercise}[$\star$]\label{exo:b1:price-theory-at-work:distance}
Shipping grain costs 0.05 dalens a tonne a kilometre. Predict the largest
price gap between two markets 300 km apart, and the gap if grain flows
between them.
\end{exercise}

\begin{exercise}[$\star\star$]\label{exo:b1:price-theory-at-work:ceiling}
For demand $Q = 80 - p$ and supply $Q = p$, compute the shortage and the
\omterm{def:b1:partial-equilibrium-welfare:dwl}{deadweight loss} of a ceiling of 20, first with costless rationing to the
buyers who value the good most, then with rationing by queues.
\end{exercise}

\begin{exercise}[$\star\star$]\label{exo:b1:price-theory-at-work:tariff}
Demand is $Q = 120 - 5p$, domestic supply $Q = 2p$, the world price 10. Compute
imports with free trade and with a tariff of 2, and show that a quota of 36
gives the same price. Compare the surpluses under the two.
\end{exercise}

\begin{exercise}[$\star\star$]\label{exo:b1:price-theory-at-work:path}
Storage costs $1\%$ of the value a month and the interest rate is $0.5\%$ a
month. A grain costs 100 at harvest. Compute its price over the next six
months under the \omterm{def:b1:price-theory-at-work:arbitrage}{arbitrage condition}.
\end{exercise}

\begin{exercise}[$\star\star$]\label{exo:b1:price-theory-at-work:rate-fall}
With depreciation at $2\%$ and no expected growth, the real rate falls from
$4\%$ to $2\%$. Predict the price-to-rent ratio before and after, and the change
in the price of a dwelling whose rent does not change.
\end{exercise}

\begin{exercise}[$\star\star$]\label{exo:b1:price-theory-at-work:corr}
The correlation of the yield with the real rate over 1950--2020 is $0.63$ in
the United States, $0.12$ in the United Kingdom and $-0.15$ in France.
% figdata: rent_corr
What do the numbers say about the rent--price condition, and which term of
it could account for the difference?
\end{exercise}

\begin{exercise}[$\star\star\star$]\label{exo:b1:price-theory-at-work:derive}
Derive the rent--price condition from the absence of a profit in buying a
dwelling, renting it out for a year and selling it at the end of the year.
\end{exercise}

\begin{exercise}[$\star\star\star$]\label{exo:b1:price-theory-at-work:random}
In \cref{exo:b1:price-theory-at-work:ceiling}, the 20 units go to buyers
drawn at random from the 60 who want one at the ceiling. Compute the value
of the units to their buyers and show that the loss exceeds the triangle.
\end{exercise}

\begin{exercise}[$\star\star\star$]\label{exo:b1:price-theory-at-work:refute}
What observation in the opening chart, or in data like it, would refute
the claim that rents equal the cost of owning?
\end{exercise}

\section{Problem: The Grain Store of Brenmoor}\label{sec:b1:price-theory-at-work:problem}

\begin{problem}[{Weekend problem --- the lean-month premium}]\label{pb:b1:price-theory-at-work:store}
Brenmoor's grain is harvested in September (invented). At harvest the
price is 200 dalens a tonne. The interest rate is $0.5\%$ a month and storing a
tonne costs 2 dalens a month. In July, demand for grain in Brenmoor is
$Q = 400 - P$, in thousands of tonnes.

\medskip\noindent\textbf{Part I --- Storage.}
\begin{enumerate}
  \item Under the \omterm{def:b1:price-theory-at-work:arbitrage}{arbitrage condition}, what is the price in October?
  \item Compute the price month by month up to July.
  \item By what percentage does the July price exceed the harvest price?
  \item Explain why the price can rise neither faster nor slower while stocks are held.
  \item If storage cost $1\%$ of the value a month instead of 2 dalens, what would the July premium be?
  \item In fact the October price is 202. What \omterm{def:b1:price-theory-at-work:carry}{convenience yield} does this reveal, per tonne for the month?
  \item What happens to the price at the next harvest, and to stocks just before it?
\end{enumerate}

\medskip\noindent\textbf{Part II --- Trade with Kestria.} Brenland's import
demand for grain is $2(220 - p_B)$ and Kestria's export supply is
$2(p_K - 160)$, in thousands of tonnes; prices in dalens a tonne.
\begin{enumerate}[resume]
  \item What are the two prices without trade?
  \item With a transport cost of 10 a tonne, find both prices and the quantity traded.
  \item Check the \omterm{def:b1:price-theory-at-work:arbitrage}{arbitrage condition}.
  \item Repeat with a transport cost of 30.
  \item At what transport cost does trade stop?
  \item Who in Kestria gains from the trade, and who loses?
\end{enumerate}

\medskip\noindent\textbf{Part III --- A ceiling in July.}
\begin{enumerate}[resume]
  \item What quantity is sold in July at the market price of Part I?
  \item The government caps the July price at 210. With the same stock, find the shortage.
  \item What queuing cost per tonne clears the market, and what is the \omterm{def:b1:price-theory-at-work:full-price}{full price}?
  \item Buyers value their time at 10 dalens an hour, and each tonne needs its own wait. How long is the wait per tonne?
  \item How much do the holders of stocks lose in July, compared with no ceiling?
  \item Predict the effect of the ceiling on storage if it is expected every year.
\end{enumerate}

\medskip\noindent\textbf{Part IV --- The store's value.} The store holds
10,000 tonnes and is rented out at 0.8 dalens a tonne a month.
\begin{enumerate}[resume]
  \item Compute the yearly rent.
  \item At a real interest rate of $5\%$, with no depreciation and no growth, what is the store worth?
  \item What would it be worth at $4\%$ and at $6\%$?
  \item What is its \omterm{def:b1:price-theory-at-work:rental-yield}{rental yield} at $5\%$?
  \item If the July ceiling of Part III were expected every year, what would happen to the store's rent and value?
  \item State the result: the lean-month price premium the store earns.
\end{enumerate}
\end{problem}
